% Generated by simulation/revision_experiments.py from declared scenarios and report.

The initial conditions and baseline input are \begin{gather*}\bm{x}(0)=(0.55,0.35,0.60,0.45,0.60,0.65,0.40),\\ \bm{u}_0=(0.20,0.35,0.30).\end{gather*}

During $5\leq t<10$, two unit-weight demand edges in layers 2 and 3 increase their marks from $0.15$ to $1$. Their directions are both $+1$ in the aligned case and $+1,-1$ during the opposed pulse. The demand-channel coupling has $\Gamma_{23}=0.2, \Gamma_{32}=0.2$ and other entries zero. All three layer scales, the support scale, and $\Psi_0$ equal $1$; one unit-weight, unit-mark support edge is constant. Thus $d_{\mathrm{ext}}$ rises from $\frac{9}{34}\approx 0.26$ to $\frac{12}{17}\approx 0.71$ in both cases, while the opposed pulse has $\chi=\frac{64}{89}\approx 0.72$ and the aligned case has $\chi=0$. The bridge coefficients are $(a_{\mathrm{conflict}},\lambda_d,\lambda_h)=(0.30,0.20,0.15)$.

The counter-mechanism uses $(k_{\mathrm{visibility}},k_{\mathrm{access}},k_{\mathrm{input}},k_{\mathrm{isolation}})=(0.10,0.08,0.20,0.20)$; the baseline sets all four to zero. The stronger-input comparison changes the input to $(0.60,0.70,0.60)$ during the same pulse. Baseline and counter-mechanism coefficients, schedules, outcomes, and solver settings are disclosed in the companion scenario files.

\begin{table}[H]\centering\small

\begin{tabularx}{\textwidth}{@{}Yrrrr@{}}\toprule

Declared synthetic comparison & $r(10)$ & $r(20)$ & $w_1(20)$ & $w_2(20)$ \\\midrule

Aligned demand & 0.5652 & 0.4307 & 0.4686 & 0.8778 \\

Opposed demand & 0.5145 & 0.4209 & 0.4524 & 0.8797 \\

Conflict pathway off & 0.5652 & 0.4307 & 0.4686 & 0.8778 \\

External-demand pathway off & 0.5529 & 0.4659 & 0.5150 & 0.7901 \\

Support pathway off & 0.5008 & 0.3832 & 0.2178 & 0.8800 \\

Two-state, opposed & 0.5808 & 0.6109 & 0.6500 & 0.8025 \\

Two-state, aligned & 0.6281 & 0.6132 & 0.6500 & 0.8025 \\

Counter-mechanism & 0.5117 & 0.4737 & 0.3513 & 0.8798 \\

Stronger inputs, zero costs & 0.5577 & 0.4634 & 0.5170 & 0.8537 \\

Stronger inputs, counter-mechanism & 0.3769 & 0.4892 & 0.3894 & 0.8555 \\

\bottomrule\end{tabularx}

\caption{Normalized model coordinates under the declared synthetic assumptions. The aligned and opposed conflict-ablation results coincide; one duplicate row is omitted.}\label{tab:synthetic-results}\end{table}

The aligned and opposed trajectories coincide when $a_{\mathrm{conflict}}=0$ (maximum state difference $0$), isolating the encoded directional pathway; those curves overlap in Figure~\ref{fig:synthetic}B. Stronger inputs increase $r(10)$ under zero costs, but decrease it under the counter-mechanism. At $t=20$ the latter comparison reverses again: the counter-mechanism gives $r=0.4737$ with baseline inputs and $r=0.4892$ with the stronger-input pulse. The sign therefore depends on outcome time and the selected cost structure. Removing support reduces the close-bond coordinate; removing external transmission reduces experienced load. These calculations identify consequences of the encoded channels.

The two-state benchmark examines direct resource/load pathways at fixed background coordinates. At $t=10$, it gives $r=0.6281$ for aligned and $0.5808$ for opposed demands. The full model adds evolving cross-coordinate dynamics, including return paths such as $r\to s\to r$ and $r\to w_1\to r$, and time-dependent feedforward modulation. The maximum resource difference between the opposed full model and its two-state comparator is $0.1900$. The comparison bundles these changes. Attributing the contrast to an individual feedback path would require targeted ablations; evaluating its empirical value requires prospective comparison.

Classical RK4 uses step $0.02$ and horizon $20$, splitting at every scheduled discontinuity. Across all 11 comparisons, the maximum state difference on the common grid between steps $0.02$ and $0.01$ is $1.46\times10^{-11}$; refinement ratios against the $0.005$ run range from $15.32$ to $16.19$. These solver checks concern the declared scenarios only.

One-factor structural sweeps use zero, half, nominal, and twice-nominal pathway coefficients. For each declared value the companion sensitivity table reports $r(10)$, $r(20)$, $w_1(20)$, and $w_2(20)$, including adverse and horizon-dependent comparisons. The coefficients remain uncalibrated throughout these structural range checks.
